\documentclass[letterpaper,landscape]{article}
\usepackage[margin=0.32in]{geometry}
\usepackage[T1]{fontenc}
\usepackage{lmodern,amsmath,amssymb,multicol,xcolor,enumitem,microtype}
\usepackage[hidelinks]{hyperref}
\hypersetup{pdftitle={STA 431 Midterm II: Calculator-priority Notesheet},pdfauthor={Prepared for Ibrahim Shahid}}
\definecolor{navy}{RGB}{24,49,72}
\definecolor{rulegray}{RGB}{168,177,184}
\pagestyle{empty}
\setlength{\parindent}{0pt}\setlength{\parskip}{1.4pt}
\setlength{\columnsep}{0.19in}\setlength{\columnseprule}{0.3pt}
\def\columnseprulecolor{\color{rulegray}}
\setlength{\abovedisplayskip}{2pt}\setlength{\belowdisplayskip}{2pt}
\setlength{\abovedisplayshortskip}{2pt}\setlength{\belowdisplayshortskip}{2pt}
\DeclareMathOperator{\E}{E}\DeclareMathOperator{\Var}{Var}\DeclareMathOperator{\Cov}{Cov}
\newcommand{\dd}{\,d}
\newcommand{\head}[1]{\par\vspace{3pt}{\color{navy}\bfseries\fontsize{10}{11}\selectfont #1}\par\vspace{1pt}\hrule\vspace{3pt}}
\newcommand{\mini}[1]{{\bfseries #1}}
\newcommand{\cmd}[1]{\par\noindent{\ttfamily\fontsize{8}{9.1}\selectfont\detokenize{#1}}\par}
\newcommand{\topline}[2]{{\color{navy}\bfseries\fontsize{15}{16}\selectfont STA 431\quad Midterm II}\hfill{\bfseries #1}\par\vspace{2pt}{\fontsize{8}{9}\selectfont #2}\par\vspace{3pt}\hrule\vspace{3pt}}
\begin{document}
\fontsize{8.3}{9.15}\selectfont
\topline{FRONT: FORMULAS + SETUP}{Handwritten-copy master\quad\textbullet\quad US Letter landscape\quad\textbullet\quad Back: native calculator workflows + STAM2 menu map}
\begin{multicols*}{3}\raggedcolumns
\head{1\quad Joint laws, support, probabilities}
\mini{Normalize.} PMF $p\ge0$, $\sum\sum p=1$; PDF $f\ge0$, $\iint_Sf=1$. If $f=cq$, then $c=1/\iint_Sq$ (sum for PMF). State $S$; zero elsewhere.
\[
 p_X(x)=\sum_yp(x,y),\quad
 f_X(x)=\int_{y:(x,y)\in S}f(x,y)\dd y.
\]
Eliminate the other variable; the marginal support is the projection of $S$. Swap $x,y$ for $Y$.
\[
 p_{X\mid Y}(x\mid y)=\frac{p(x,y)}{p_Y(y)},\quad
 f_{X\mid Y}(x\mid y)=\frac{f(x,y)}{f_Y(y)}.
\]
Denominator positive; keep the slice of $S$ at the given $y$; normalize over $x$.
$f(x,y)=f_{X\mid Y}(x\mid y)f_Y(y)$.
Two marginals alone, or one conditional alone, generally do not determine the joint.

\mini{CDF.} $F(x,y)=P(X\le x,Y\le y)$.
\[
 F(x,y)=\sum_{u\le x}\sum_{v\le y}p(u,v)
 \quad\text{or}\quad\int_{-\infty}^{x}\int_{-\infty}^{y}f(u,v)\dd v\dd u
\]
(use the applicable sum or integral). Intersect with $S$.
\[
 F_X(x)=\lim_{y\to\infty}F(x,y),\quad
 f(x,y)=\frac{\partial^2F}{\partial x\partial y}\ \text{a.e.}
\]
The derivative formula needs an absolutely continuous law.
\[
\begin{split}
P(a<X\le b,c<Y\le d)={}&F(b,d)-F(a,d)\\
&-F(b,c)+F(a,c).
\end{split}
\]
For integer inclusive lower endpoints use $a-1,c-1$; otherwise left limits. Integer PMF = $F(x,y)-F(x-1,y)-F(x,y-1)+F(x-1,y-1)$.

\mini{Event bounds.} $P((X,Y)\in A)=\iint_{S\cap A}f$ or sum allowed pairs. Draw the intersection and split where a boundary changes. Inner limits may depend on the outer variable.
On $0<x<1,0<y<2$,
\[
 P(X<Y)=\int_0^1\int_x^2f\dd y\dd x.
\]
On $x,y>0,x+y<1$, use $0<x<1$, $0<y<1-x$. For $X+Y\le a$, intersect $y\le a-x$ with $S$.

\mini{Words.} Both survive $t$: $X>t,Y>t$. Series fails if either fails: $T=\min(X,Y)$,
$P(T\le t)=1-P(X>t,Y>t)$. Parallel fails if both fail: $T=\max(X,Y)$, $P(T\le t)=F(t,t)$.

\mini{Independence.} $F(x,y)=F_X(x)F_Y(y)$ for all $x,y$; equivalent to $p=p_Xp_Y$ at every pair, or $f=f_Xf_Y$ almost everywhere, including zero regions. A point failure disproves discrete independence; use positive-area failure for PDFs. Product formula on triangular support is insufficient.
\par\columnbreak\par
\head{2\quad Expectation, moments, MGFs}
\mini{LOTUS.}
\[
 \E[g(X)]=\sum_xg(x)p_X(x)\quad\text{or}\quad\int g(x)f_X(x)\dd x.
\]
For $g(X,Y)$ use the joint sum/integral. A finite expectation needs $\E|g|<\infty$.
\[
 \E(aX+bY+c)=a\E X+b\E Y+c.
\]
Linearity needs no independence. Generally $\E g(X)\ne g(\E X)$; $\E1_A=P(A)$.

\mini{Raw and central moments.}
\[
 \mu_r'=\E X^r,\quad\mu=\mu_1',\quad
 \mu_r=\E[(X-\mu)^r].
\]
\[
 \mu_0'=1,\quad\mu_1=0,\quad
 \Var X=\mu_2=\mu_2'-\mu^2.
\]
\[
\begin{aligned}
\mu_3&=\mu_3'-3\mu\mu_2'+2\mu^3,\\
\mu_4&=\mu_4'-4\mu\mu_3'+6\mu^2\mu_2'-3\mu^4.
\end{aligned}
\]
\[
 \Var(aX+b)=a^2\Var X,\quad
 \sigma_{aX+b}=|a|\sigma_X.
\]
$Z=(X-\mu)/\sigma$ has mean $0$, variance $1$; not necessarily normal. Symmetry about zero makes existing odd moments vanish.
For $Y=g(X)$, $\Var Y=\E[g(X)^2]-(\E g(X))^2$; $Y=X^2$ needs $\E X^4$.

\mini{MGF definition.}
\[
 M_X(t)=\E e^{tX}=\sum_xe^{tx}p_X(x)
 \quad\text{or}\quad\int e^{tx}f_X(x)\dd x.
\]
State the convergence domain. Standard moment/uniqueness theorems need finiteness on an open interval containing $0$.
\[
\begin{gathered}
 M(0)=1,\quad M^{(r)}(0)=\E X^r,\\
 \Var X=M''(0)-[M'(0)]^2.
\end{gathered}
\]
The coefficient of $t^r$ is $\E X^r/r!$.
For $K(t)=\ln M(t)$, $K'(0)=\mu$, $K''(0)=\sigma^2$.
Equal MGFs near zero imply equal distributions; equal means/variances do not.
\[
 M_{aX+b}(t)=e^{bt}M_X(at).
\]
For independent $X_i$,
\[
 M_{\sum a_iX_i+b}(t)=e^{bt}\prod_iM_{X_i}(a_it).
\]
Intersect transformed domains; without independence do not multiply marginal MGFs.

\mini{Fast forms.} Rate-$\lambda$ exponential: $M(t)=\lambda/(\lambda-t)$ for $t<\lambda$; mean $1/\lambda$, variance $1/\lambda^2$.
If $M(t)=e^{\alpha t+\beta t^2}$, mean $\alpha$, variance $2\beta$ ($\beta\ge0$).
$\int_0^\infty x^ne^{-ax}\dd x=n!/a^{n+1}$ for $a>0$, integer $n\ge0$.
\par\columnbreak\par
\head{3\quad Products, combinations, conditioning}
\mini{Product moments.}
\[
 \mu'_{rs}=\E[X^rY^s],\quad
 \mu_{rs}=\E[(X-\mu_X)^r(Y-\mu_Y)^s].
\]
\[
 \Cov(X,Y)=\E XY-\mu_X\mu_Y,\quad
 \rho=\frac{\Cov(X,Y)}{\sigma_X\sigma_Y}.
\]
Correlation needs positive finite variances; $|\rho|\le1$ and $|\Cov(X,Y)|\le\sigma_X\sigma_Y$.
\[
\begin{gathered}
 \Cov(X,Y)=\Cov(Y,X),\quad\Cov(X,X)=\Var X,\\
 \Cov(aX+b,cY+d)=ac\Cov(X,Y).
\end{gathered}
\]
Independence implies zero covariance (finite second moments); converse false.
For $X\sim U(-1,1)$, $Y=X^2$, $\E X=\E X^3=0$ so covariance is zero. Yet $A=\{|X|<1/2\}$, $B=\{Y<1/4\}$ give $P(A\cap B)=1/2\ne1/4=P(A)P(B)$.

\mini{One combination.} $W=c+\sum_i a_iX_i$:
\[
\begin{aligned}
 \E W&=c+\sum_i a_i\mu_i,\\
 \Var W&=\sum_i a_i^2\sigma_i^2
 +2\sum_{i<j}a_ia_j\sigma_{ij}.
\end{aligned}
\]
Here $\sigma_{ij}=\Cov(X_i,X_j)$. Constants vanish in variances; preserve coefficient signs in cross terms; each unordered pair appears once, times $2$. Zero covariance suffices to remove its cross term.

\mini{Two combinations.} $U=c+\sum_i a_iX_i$, $V=d+\sum_i b_iX_i$:
\[
 \Cov(U,V)=\sum_i\sum_j a_ib_j\sigma_{ij}.
\]
Compute both variances for $\rho_{UV}$. A sum over all $i,j$ already counts both off-diagonal orders; do not insert another factor $2$. Matrix shortcut is on the back.

\mini{Conditional moments.} Fix the given $y$:
\[
 m(y)=\E[X\mid Y=y]=\frac{\int x f(x,y)\dd x}{f_Y(y)}
\]
(use sums and $p_Y$ for discrete). Replace $x$ by $g(x)$ for $\E[g(X)\mid Y=y]$. Integrate only the conditional slice; denominator positive.
\[
 \Var(X\mid Y=y)=\E[X^2\mid Y=y]-m(y)^2.
\]
$m(y_0)$ is a number; $\E[X\mid Y]=m(Y)$ is a random variable.
\[
\begin{gathered}
 \E X=\E[\E(X\mid Y)],\\
 \Var X=\E[\Var(X\mid Y)]+\Var(\E[X\mid Y]).
\end{gathered}
\]
Independence gives $\E[g(X)\mid Y]=\E g(X)$; a constant conditional mean alone does not prove independence. For an event $B$, $P(A\mid B)=P(A\cap B)/P(B)$. Continuous $Y=y$ uses a density, not division by $P(Y=y)=0$.
\end{multicols*}
\newpage
\topline{BACK: TI-84 PLUS CE WORKFLOWS}{\texttt{PRGM > STAM2}: 1 Exam assistant; 2 General solvers; 3 Encyclopedia; 4 Practice; 5 Help; 6 Quit. Enter your current data; -> means STO\,$\blacktriangleright$}
\begin{multicols*}{3}\raggedcolumns
\head{4\quad Lists: most discrete calculations}
\mini{1-Var Stats: let weights do the work.}
Enter values in $L_1$, nonnegative weights in $L_2$; at least one weight positive. Weights need not sum to $1$.
\cmd{1-Var Stats L1,L2}
\texttt{STAT > CALC}; wizard: List $L_1$, FreqList $L_2$. Read $\bar x=\E X$, $\sigma_x^2=\Var X$, $n=\sum w$. Use population $\sigma_x$, not sample $S_x$; fractional weights can leave $S_x$ undefined.

HW5 Q1: $L_1=\{1,2,3,4\}$, $L_2=\{2,6,12,20\}$. No need to divide the list by $40$ first. Output $\bar x=13/4$, $\sigma_x^2=63/80$, $c=1/n=1/40$.

\mini{Functions of $X$.} Put transformed values in another list, keeping weights:
\cmd{L1^2->L3}
\cmd{1-Var Stats L3,L2}
Read $\E[X^2]$ and $\Var(X^2)$ directly. For any $g$, use $g(L_1)$ in $L_3$. For a raw moment,
\cmd{sum(L1^R*L2)/sum(L2)}
For a central moment, store the mean in $M$, then use \texttt{sum((L1-M)\^{}R*L2)/sum(L2)}.

\mini{Joint PMF: one row per allowed pair.}
$L_1$: every $x$; $L_2$: matching $y$; $L_3$: nonnegative joint weights. Repeat coordinates as needed. Let
\cmd{sum(L3)->S}
If already probabilities, $S=1$. Otherwise normalizing constant is $1/S$. General product moment:
\cmd{sum(L1^R*L2^T*L3)/S}
For example $R=T=1$ gives $\E XY$.

\mini{Events: masks select rows.}
\cmd{sum((L1>L2)*L3)/S}
\cmd{sum((L1+L2<=3)*L3)/S}
These give $P(X>Y)$ and $P(X+Y\le3)$. Use the single $\le$ token for printed \texttt{<=}, from \texttt{[2nd][MATH] TEST}.

\mini{Conditionals: mask then reuse Stats.}
For $X\mid Y=A$, store the given value in $A$:
\cmd{L3*(L2=A)->L4}
\cmd{1-Var Stats L1,L4}
Then $\bar x$ is the conditional mean and $\sigma_x^2$ the conditional variance. The mask's total must be positive. For a marginal mass,
\cmd{sum((L2=A)*L3)/S}
For $Y\mid X=A$, mask with \texttt{L1=A} and run Stats on $L_2$. For $P(X>H\mid Y=A)$ divide \texttt{sum((L1>H)*L4)} by \texttt{sum(L4)}.
\par\columnbreak\par
\head{5\quad Continuous: one reusable setup}
\mini{Save the shape in $Y_1$.} In \texttt{Y=}, enter nonnegative unnormalized $q(X)$; bounds $B,C$. For a conditional slice $q(x,y)$, let $X$ be the integration variable and $A$ the fixed given value; $B,C$ must match that slice. Example HW5 Q7: $Y_1=X+A$, $A=1/2$, $B=0$, $C=1$.
\cmd{fnInt(Y1,X,B,C)->D}
\cmd{fnInt(X*Y1,X,B,C)/D->M}
\cmd{fnInt((X-M)^2*Y1,X,B,C)/D}
These give mass $D$, mean $M$, and variance. Normalization cancels in the ratios. For a full unnormalized density, $c=1/D$; for a joint slice, $D$ is the appropriate marginal value if the joint is already normalized.

\mini{Any moment or payoff.}
\cmd{fnInt(X^R*Y1,X,B,C)/D}
Replace \texttt{X\^{}R} by $g(X)$ for $\E g(X)$. To get $\Var(X^2)$, first store
\cmd{fnInt(X^2*Y1,X,B,C)/D->M}
then evaluate
\cmd{fnInt((X^2-M)^2*Y1,X,B,C)/D}
The centered form avoids subtracting two nearly equal large moments. HW5 Q2: $Y_1=1-X^2$, $B=-1$, $C=1$ gives $\Var X=1/5$, $\E X^2=1/5$, $\Var(X^2)=8/175$.

\mini{Probability by ratio.} For $[E,F]$ within $[B,C]$,
\cmd{fnInt(Y1,X,E,F)/D}
For $P(X>E)$, integrate from $E$ to $C$; clip bounds to the support. For HW4 Q2 conditional $X\mid Y=A$, use $Y_1=2X+A$, $B=0,C=1$. For each new $A$, recompute $D,M$.

\mini{Double integrals.} App: \texttt{GENERAL SOLVERS > DENSITY > DOUBLE}. Enter $q(X,Y)$, outer $x$ limits, inner $y$ limits, and $g=1$ for probability. Use RAW integral for a normalized PDF on an event. Compare $N,2N$. Native \texttt{fnInt} cannot nest; reduce one integral by hand. HW4 Q5 gives
\cmd{fnInt(1/2+X-5*X^2/8,X,0,1)}
which gives $19/24$. Split a changing boundary before calculating. A one-dimensional marginal at a fixed value can use \texttt{fnInt} directly; a symbolic marginal formula still needs algebra.

\mini{MGF value from a PDF.} With $T$ a specified number,
\cmd{fnInt(e^(T*X)*Y1,X,B,C)/D}
computes $M(T)$ on a finite support. It does not establish convergence on an infinite support or provide a symbolic MGF. Avoid using arbitrary large finite bounds as if they were $\infty$.
\par\columnbreak\par
\head{6\quad Covariance, matrices, MGFs, keys}
\mini{2-Var Stats from joint rows.}
Using $L_1=x,L_2=y,L_3=w$ from column 4,
\cmd{2-Var Stats L1,L2,L3}
read $\bar x,\bar y,\sigma_x,\sigma_y,n,\Sigma xy$:
\[
 \Cov(X,Y)=\frac{\Sigma xy}{n}-\bar x\bar y,
 \quad\rho=\frac{\Cov(X,Y)}{\sigma_x\sigma_y}.
\]
Get actual statistics tokens through \texttt{VARS > Statistics}. Store results before editing lists or running another Stats command. Direct list alternative:
\cmd{sum(L1*L3)/S->M}
\cmd{sum(L2*L3)/S->N}
\cmd{sum((L1-M)*(L2-N)*L3)/S}
This is covariance, with $S=\texttt{sum(L3)}$.

\mini{Matrices: automate every cross term.}
\texttt{[2nd][}$x^{-1}$\texttt{] MATRIX > EDIT}. Enter covariance matrix $[A]=\Sigma$ with variances on its diagonal. Put coefficients of one or more combinations in rows of $[B]$. Then
\[
 [B][A][B]^{\mathsf T}
\]
is their covariance matrix. One row gives a $1\times1$ variance. Two rows give $\begin{pmatrix}\Var U&\Cov(U,V)\\\Cov(U,V)&\Var V\end{pmatrix}$. Ignore additive constants. The transpose is the actual postfix $\mathsf T$ token from \texttt{MATRIX > MATH}, not alpha-T.
HW5 Q6:
\[
 [A]=\begin{pmatrix}4&1&-1\\1&9&0\\-1&0&2\end{pmatrix},\quad
 [B]=\begin{pmatrix}3&-2&1\end{pmatrix}.
\]
Type $[B]*[A]*[B]^{\mathsf T}$ to obtain $[[56]]$. If given correlation, first use $\Cov(X,Y)=\rho\sigma_X\sigma_Y$.

\mini{A given MGF in $Y_1$.} Enter $M_X(X)$; $X$ now represents $t$. For $Z=aX+b$, store $a$ in $A$, $b$ in $B$; put this transform in $Y_2$:
\cmd{e^(B*X)*Y1(A*X)}
Use actual $Y_1$ token from \texttt{VARS > Y-VARS > Function}. The mean is checked by
\cmd{nDeriv(Y1,X,0)}
App \texttt{MGF TOOLS} estimates first/second derivatives at two step sizes. For exact work recognize $e^{\alpha t+\beta t^2}$: mean $\alpha$, variance $2\beta$. Keep stable digits; numerical derivatives are approximations.

\mini{Fast entry and reliable use.} \texttt{[2nd][0]}: CATALOG. \texttt{[2nd][STAT]}: LIST; \texttt{seq} under OPS, \texttt{sum} under MATH. A finite sum is
\cmd{sum(seq(expression,K,first,last))}
Store reusable numbers with STO\,$\blacktriangleright$; do not rely on \texttt{Ans} after another command. \texttt{MATH > Frac} suggests a fraction, not a proof. Calculator output handles arithmetic; write the support, event, formula, bounds, MGF domain, and independence justification yourself.
\end{multicols*}
\end{document}
